\section*{Hoofdstuk \ref{ch:g7:lungs-and-gas-exchange} --- De longen en de gaswisseling}

\begin{solution}{exo:g7:lungs-and-gas-exchange:1}
\omterm{def:g4:breathing:movements}{Inademen}: de ribspieren tillen de kast op en maken haar breder terwijl het
\omterm{prop:g7:lungs-and-gas-exchange:ventilation}{middenrif} naar beneden afvlakt; de opgerekte borst trekt de \omterm{def:g7:how-animals-breathe:four}{longen} wijd en
er stroomt lucht naar binnen --- spierwerk. \omterm{def:g4:breathing:movements}{Uitademen} in rust: de \omterm{def:g3:bones-and-muscles:muscle}{spieren}
ontspannen en de opgerekte borst veert terug en duwt de lucht naar buiten
--- loslaten, vrijwel gratis.
\end{solution}

\begin{solution}{exo:g7:lungs-and-gas-exchange:2}
De brede spiervloer onder de \omterm{def:g7:how-animals-breathe:four}{longen}, die de borst van onderen afsluit; het
afvlakken ervan is de diepe helft van elke ademteug naar binnen.
\end{solution}

\begin{solution}{exo:g7:lungs-and-gas-exchange:3}
De luchtzakjes van de \omterm{def:g7:how-animals-breathe:four}{longen}: honderden miljoenen dunwandige bolletjes aan
de fijnste uiteinden van de luchtwegen, elk omwikkeld met fijne
\omterm{def:g4:heart-and-blood:vessels}{bloedvaten}.
\end{solution}

\begin{solution}{exo:g7:lungs-and-gas-exchange:4}
Ingeademd: ongeveer $21$ procent \omterm{def:g7:what-is-respiration:respiration}{zuurstof}, vrijwel geen \omterm{def:g7:what-is-respiration:respiration}{koolstofdioxide}.
Uitgeademd: ongeveer $16$ procent \omterm{def:g7:what-is-respiration:respiration}{zuurstof} en ongeveer $4$ procent
\omterm{def:g7:what-is-respiration:respiration}{koolstofdioxide}. De verschillen --- vijf punten \omterm{def:g7:what-is-respiration:respiration}{zuurstof} weggenomen, vier
punten \omterm{def:g7:what-is-respiration:respiration}{koolstofdioxide} toegevoegd --- zijn de wisseling zelf.
\end{solution}

\begin{solution}{exo:g7:lungs-and-gas-exchange:5}
Het groot-dun-vochtige, met \omterm{prop:g4:journey-of-food:absorb}{bloed} doorregen oversteekoppervlak. Gedeeld
met de \omterm{def:g1:animals-around-us:covering}{veren} van de \omterm{def:g7:how-animals-breathe:four}{kieuwen}
(\cref{prop:g7:how-animals-breathe:spec}) en de \omterm{prop:g7:digestion-and-nutrients:absorption}{darmvlokken}
(\cref{ex:g7:digestion-and-nutrients:spec}).
\end{solution}

\begin{solution}{exo:g7:lungs-and-gas-exchange:6}
Teer die de luchtwegen beslaat en prikkelt (de overbelaste bezem en zijn
hoest); wanden van \omterm{def:g7:lungs-and-gas-exchange:alveoli}{longblaasjes} die afbreken --- oppervlak verloren; een
rookgas dat de zuurstofstoelen van het \omterm{prop:g4:journey-of-food:absorb}{bloed} bezet; nauwer geworden
luchtwegen die de blaasbalg smoren.
\end{solution}

\begin{solution}{exo:g7:lungs-and-gas-exchange:7}
Naar binnen moet borst en \omterm{def:g7:how-animals-breathe:four}{longen} tegen hun \omterm{def:g1:animals-around-us:covering}{veer} in oprekken --- dat
oprekken is het werk. Het \omterm{def:g4:breathing:movements}{uitademen} in rust is die \omterm{def:g1:animals-around-us:covering}{veer} die terugkeert: de
verbrede borst en de opgerekte \omterm{def:g7:how-animals-breathe:four}{longen} \omterm{def:g1:animals-around-us:covering}{veren} terug, en dat terugveren duwt
de lucht.
\end{solution}

\begin{solution}{exo:g7:lungs-and-gas-exchange:8}
De wisseling heeft net zo goed oppervlak nodig als verse lucht: de
oversteek gebeurt alleen waar lucht een met \omterm{prop:g4:journey-of-food:absorb}{bloed} doorregen wand raakt.
Een snellere blaasbalg door minder, slappere blaasjes ventileert oppervlak
dat niet meer bestaat --- de ontbrekende factor is het oversteekoppervlak
zelf.
\end{solution}

\begin{solution}{exo:g7:lungs-and-gas-exchange:9}
De roker klimt met nauwere luchtwegen (gesmoorde inlaat), een verkleind
oppervlak aan \omterm{def:g7:lungs-and-gas-exchange:alveoli}{longblaasjes} (minder lading per ademteug) en zuurstofstoelen
die door het rookgas bezet zijn (minder vervoerd per ronde). Om dezelfde
spierrekening te betalen moet het \omterm{def:g4:heart-and-blood:heart}{hart} meer rondes draaien --- vandaar de
hijgende aankomst bij een \omterm{prop:g4:heart-and-blood:pulse}{polsslag} op sprintniveau.
\end{solution}

\begin{solution}{exo:g7:lungs-and-gas-exchange:10}
Allebei zijn het getrainde beheersing van de \emph{uitgaande} fase: er
wordt afgemeten spierkracht toegevoegd aan wat normaal een passieve \omterm{def:g1:animals-around-us:covering}{veer}
is --- de duiker om uitzettende lucht gestaag af te blazen, de speler om
één lange, beheerste uitstroom te doseren. (Hun diepe teugen trainen ook
het werk naar binnen.)
\end{solution}

\begin{solution}{exo:g7:lungs-and-gas-exchange:11}
Een diepere en snellere \omterm{def:g7:what-is-respiration:respiration}{ademhaling} vermenigvuldigt het laden aan het
loket van de \omterm{def:g7:lungs-and-gas-exchange:alveoli}{longblaasjes}; het snellere, hardere \omterm{def:g4:heart-and-blood:heart}{hart} vermenigvuldigt de
rondes aan het transportloket; het omleiden verwijdt de plaatselijke
vaten van de \omterm{def:g3:bones-and-muscles:muscle}{spieren} aan het bezorgloket. Drie \omterm{def:g4:adapted-to-their-world:adaptation}{aanpassingen}, één
vermenigvuldigde levering (de kloppende boeken van
\cref{ex:g7:lungs-and-gas-exchange:accounting}).
\end{solution}

\begin{solution}{exo:g7:lungs-and-gas-exchange:12}
De grens: honderden miljoenen \omterm{def:g7:lungs-and-gas-exchange:alveoli}{longblaasjes} tellen op tot ongeveer een
halve tennisbaan uitwisselingsoppervlak, opgevouwen in een borstkas ---
en daaroverheen steken vijf procentpunten van de \omterm{def:g7:what-is-respiration:respiration}{zuurstof} van elke
ademteug naar binnen over en vier punten \omterm{def:g7:what-is-respiration:respiration}{koolstofdioxide} naar buiten. Het
trage vuur: het teer en de prikkeling van rook verteren die grens jaar na
jaar --- de aanlooproutes beslagen, de blaasjes ingestort, de
transportstoelen bezet --- cumulatief, stil en maar deels omkeerbaar, als
een grens die sneller afbrandt dan hij herbouwd kan worden.
\end{solution}

\begin{solution}{pb:g7:lungs-and-gas-exchange:1}
\textbf{1.} Fase naar binnen: de ribspieren tillen de kast op en maken
haar breder, het \omterm{prop:g7:lungs-and-gas-exchange:ventilation}{middenrif} vlakt af --- dat kost \omterm{prop:g3:balanced-diet:jobs}{energie}. Fase naar
buiten (in rust): de \omterm{def:g3:bones-and-muscles:muscle}{spieren} laten los, borst en \omterm{def:g7:how-animals-breathe:four}{longen} \omterm{def:g1:animals-around-us:covering}{veren} terug ---
vrijwel gratis. De cyclus ongeveer $15$ keer per minuut.

\textbf{2.} De \omterm{def:g7:how-animals-breathe:four}{longen} kleven aan de binnenwand van de borst: verbreed de
borst en de \omterm{def:g7:how-animals-breathe:four}{longen} worden ermee wijd getrokken, en de buitenlucht stroomt
naar binnen om de opgerekte ruimte te vullen. De pomp werkt door de kamer
te vervormen, niet door het orgaan te knijpen.

\textbf{3.} Ongeveer $15$ cycli per minuut in rust; bij inspanning worden
de cycli sneller \emph{én} dieper --- meer teugen, en meer lucht per teug.

\textbf{4.} Hijgen: snelle oppervlakkige cycli, bijna helemaal \omterm{prop:g7:lungs-and-gas-exchange:ventilation}{middenrif}.
Zuchten: één extra diepe rek en een lange ontspanning --- een reset van de
\omterm{def:g1:animals-around-us:covering}{veer}. Kaarsen: de uitgaande fase met spierkracht toegevoegd aan de \omterm{def:g1:animals-around-us:covering}{veer}.

\textbf{5.} De \omterm{def:g7:lungs-and-gas-exchange:alveoli}{longblaasjes}: honderden miljoenen dunwandige bolletjes,
elk omwikkeld met fijne vaten, samen ongeveer een halve tennisbaan.

\textbf{6.} Groot: de opgetelde baan. Dun: bolwanden ter dikte van een
\omterm{def:g5:first-look-at-cells:cell}{cel}. Vochtig: hun bekleding. Met \omterm{prop:g4:journey-of-food:absorb}{bloed} doorregen: de omwikkelende vaten.
Op alle vier de punten voldoet het aan het bestek.

\textbf{7.} Vijf punten achtergehouden van $21$ --- ruwweg een kwart van
de ingeademde \omterm{def:g7:what-is-respiration:respiration}{zuurstof} wordt per teug opgenomen.

\textbf{8.} Van het vocht van het oppervlak: de vertrekkende lucht is over
natte (en warme) wanden gestreken, dus voert de uitlaat water en warmte
af --- de beslagen spiegel is de onderhoudskost van het oppervlak, bij
elke ademteug betaald.

\textbf{9.} Inlaat: teer beslaat en prikkelt de luchtwegen. Pomp: nauwere
luchtwegen smoren elke cyclus. Oppervlak: wanden van \omterm{def:g7:lungs-and-gas-exchange:alveoli}{longblaasjes} breken
af en smelten samen --- oppervlak verloren. Transport: het rookgas bezet
de zuurstofstoelen van het \omterm{prop:g4:journey-of-food:absorb}{bloed}.

\textbf{10.} Hoest en verkoudheden ontsteken en ruimen op --- de
machinerie herstelt en gaat weer in dienst. Gebroken wanden van
\omterm{def:g7:lungs-and-gas-exchange:alveoli}{longblaasjes} bouwen zichzelf niet op: elk samengesmolten blaasje is
uitwisselingsoppervlak dat definitief van de balans geschrapt is ---
kapitaal, geen lopende kosten.

\textbf{11.} (1) Laat geen rook toe: kapitaalverlies is voorgoed. (2)
Draai op schone lucht: filters en bezems hebben een eindige capaciteit.
(3) Laat de pomp onder belasting werken: diepe inspanning onderhoudt de
\omterm{def:g3:bones-and-muscles:muscle}{spieren} en ventileert het hele oppervlak. (4) Neem de inlaat door de \omterm{def:g1:the-five-senses:sense}{neus}
in kou en stof: voorbewerking beschermt het werk.

\textbf{12.} Als enige van de pompen van het lichaam gehoorzaamt de
blaasbalg zowel de automatische diensten als de wil --- omdat dezelfde
uitstroom die \omterm{def:g7:what-is-respiration:respiration}{koolstofdioxide} afblaast de wind van spraak en zang is, en
een machine die aan de taal is uitgeleend de hand van de spreker op de
\omterm{def:g6:life-through-seasons:buds}{knoppen} moet accepteren.
\end{solution}
